\begin{theorem}[LambdaCalc constructor source disjointness]
\label{thm:lambdacalc-constructor-source-disjointness}
Let a carried $\mathsf{BHist}$ term row satisfy
$\mathsf{LambdaCalcBHistTermCarrier}$, and suppose its visible $\TreeUp$ root is
one of $\mathsf{lamVar}$, $\mathsf{lamAbs}$, or $\mathsf{lamApp}$. Then the
public LambdaCalc classifier reads only the source fields of that root: a
variable row exposes only its unary variable-index field, an abstraction row
exposes only its binder-index and body fields, and an application row exposes
only its left and right subterm fields. No accepted row can simultaneously be
read as a different LambdaCalc constructor under the same classifier endpoint.
\end{theorem}
\begin{proof}
Read the visible root tag through the $\TreeUp$ row in the carried
$\mathsf{BHist}$ packet. The three LambdaCalc constructors are distinct
TreeUp tags, so the tag readback fixes which source slots are present in the
packet.

In the variable case, apply
\autoref{thm:lambdacalc-variable-constructor-carrier}. Its packet contains the
unary variable-index row and no binder-body or binary subterm slots. Transport
along the displayed $\hsame$ endpoint preserves that slot inventory.

In the abstraction case, apply
\autoref{thm:lambdacalc-abstraction-application-carrier-closure}. The source packet consists
of the binder-index row and the carried body row, and $\hsame$ transport cannot
create the variable-only or application-only source slots.

In the application case, apply
\autoref{thm:lambdacalc-abstraction-application-carrier-closure}. The packet consists of the
two carried subterm rows, with no unary variable-index-only source and no
binder-body abstraction source.

Since the classifier is read from the same $\TreeUp$ tag, the same $\hsame$
endpoint, and the same $\Pkg$ provenance in every case, two different
constructor readings would require one visible root tag to be two distinct
TreeUp constructors. The tag readback excludes this, giving constructor source
disjointness.
\end{proof}

\begin{theorem}[LambdaCalc substitution source coverage]
\label{thm:lambdacalc-substitution-source-coverage}
Every accepted LambdaCalc substitution $\Cont$ ledger endpoint is sourced from
exactly one variable, abstraction, or application branch of the displayed
$\TreeUp$ syntax spine, together with the $\NatUp$ unary index rows named by the
free-variable packet. Equivalently, the substitution endpoint has no additional
source outside the constructor branch chosen by the visible syntax spine, the
binder-index ledger, the free-variable ledger, the $\hsame$ endpoint rows, and
the $\Pkg$ provenance.
\end{theorem}
\begin{proof}
Induct on the displayed $\TreeUp$ syntax spine of the carried source term. The
base variable branch exposes the variable-index row and the matching
free-variable packet endpoint, so the substitution continuation is sourced from
that single variable branch.

For an abstraction spine, the visible constructor supplies the binder-index row
and the body subterm. The $\NatUp$ unary row records the binder index, while the
free-variable packet names the rows that survive capture-avoiding transport
through the body.

For an application spine, the visible constructor supplies the left and right
subterm rows. The induction hypotheses cover the substitution endpoints of the
two subspines, and the application packet repacks those endpoints into the
single application branch selected by the root tag.

At each step, use
\autoref{thm:lambdacalc-free-variable-ledger-exactness} to identify the
free-variable rows and
\autoref{thm:lambdacalc-substitution-carrier-preservation} to transport the
carried endpoint across the displayed $\Cont$ route.

The repacking step records the same $\hsame$ endpoint rows and the same $\Pkg$
provenance. Hence every accepted substitution endpoint is covered by exactly one
constructor branch of the displayed syntax spine and the named $\NatUp$
free-variable rows.
\end{proof}

\begin{theorem}[LambdaCalc alpha-beta ledger determinacy]
\label{thm:lambdacalc-alpha-beta-ledger-determinacy}
Let two accepted LambdaCalc packets expose the same $\TreeUp$ syntax row, the
same $\NatUp$ binder-index row, the same alpha ledger, the same beta-redex row,
the same substitution $\Cont$ ledger, and the same $\Pkg$ provenance. Then their
public LambdaCalc classifier endpoints agree by $\hsame$.
\end{theorem}
\begin{proof}
First transport the common syntax and binder rows through
\autoref{thm:lambdacalc-alpha-beta-carrier-transport}. This aligns the carried
LambdaCalc endpoints determined by the shared $\TreeUp$ syntax and $\NatUp$
binder-index rows.

Next apply
\autoref{thm:lambdacalc-substitution-output-determinacy} to the shared
substitution $\Cont$ ledger. The two packets therefore expose the same
substitution output endpoint after the displayed endpoint transports.

Finally apply
\autoref{thm:lambdacalc-alpha-beta-ledger-separation}. The common alpha ledger
and common beta-redex row cannot introduce a second public endpoint once the
syntax, binder, substitution ledger, and provenance rows have been fixed.

The remaining comparison is componentwise $\hsame$ transport along the shared
rows. Repacking those transports with the common $\Pkg$ provenance gives the
same public LambdaCalc classifier endpoint for both packets.
\end{proof}

\begin{theorem}[LambdaCalc root constructor frontier]
\label{thm:lambdacalc-root-constructor-frontier}
The public LambdaCalc term frontier exposed by a carried $\mathsf{BHist}$ row is
exhausted by the three visible root constructors
$\mathsf{lamVar}$, $\mathsf{lamAbs}$, and $\mathsf{lamApp}$. Each accepted root
row reads exactly the $\TreeUp$ source slots, $\NatUp$ index rows, $\hsame$
endpoint rows, and $\Pkg$ provenance assigned to that constructor branch, and no
fourth public root constructor can be read from the same packet.
\end{theorem}
\begin{proof}
Start with the carried $\mathsf{BHist}$ packet and read its visible $\TreeUp$
root. The constructor packet theorems cover the variable, abstraction, and
application cases as the public constructors of the LambdaCalc carrier.

Apply \autoref{thm:lambdacalc-constructor-source-disjointness}. It separates the
source slots attached to the three visible roots and rules out a second
constructor reading at the same endpoint.

For substitution or beta-visible rows attached to the root, use
\autoref{thm:lambdacalc-substitution-source-coverage}. Those rows are still
sourced from the constructor branch selected by the same syntax spine and the
named $\NatUp$, $\hsame$, and $\Pkg$ rows.

Thus the public frontier is exactly the listed constructor inventory, with all
readable rows anchored in the displayed $\mathsf{BHist}$ packet.
\end{proof}

\begin{theorem}[LambdaCalc root substitution frontier]
\label{thm:lambdacalc-root-substitution-frontier}
For every accepted LambdaCalc root substitution endpoint, the substitution
provenance is preserved through the displayed $\Cont$ continuation and the
carried $\Pkg$ rows. The endpoint remains attached to the same visible
$\mathsf{BHist}$ constructor branch, the same $\TreeUp$ syntax spine, the same
$\NatUp$ binder and free-variable rows, and the same componentwise $\hsame$
transport used by the source packet.
\end{theorem}
\begin{proof}
Begin with the accepted root substitution endpoint. Apply
\autoref{thm:lambdacalc-substitution-source-coverage} to identify the unique
constructor branch, binder-index row, free-variable packet, and $\Pkg$
provenance that source the endpoint.

Next use \autoref{thm:lambdacalc-namecert-substitution-ledger}. It reads the
endpoint from the displayed $\mathsf{BHist}$ traversal and capture-avoidance
$\Cont$ ledger, so no ambient syntax equality or unrecorded variable set is
available as a second source.

Finally apply
\autoref{thm:lambdacalc-substitution-composition-carrier-closure}. Composition
of the accepted continuations preserves the carried LambdaCalc endpoint and
reuses the same $\Pkg$ provenance. The root substitution frontier is therefore
stable under the displayed $\Cont$ and $\Pkg$ rows.
\end{proof}

\begin{theorem}[LambdaCalc root beta-ledger frontier]
\label{thm:lambdacalc-root-beta-ledger-frontier}
Suppose two accepted LambdaCalc beta-ledger root packets expose the same
$\TreeUp$ syntax row, the same $\NatUp$ binder-index row, the same alpha ledger,
the same beta-redex row, the same substitution $\Cont$ ledger, and the same
$\Pkg$ provenance. Under the displayed alpha-beta classifier, the two root
packets expose the same public LambdaCalc endpoint by componentwise $\hsame$.
\end{theorem}
\begin{proof}
First apply \autoref{thm:lambdacalc-alpha-beta-ledger-separation}. The alpha
ledger and beta-redex row occupy disjoint readable positions once the syntax
spine, binder row, substitution ledger, and provenance are fixed.

Next apply \autoref{thm:lambdacalc-alpha-beta-ledger-determinacy}. Its
hypotheses are exactly the shared $\TreeUp$, $\NatUp$, alpha, beta,
substitution, and $\Pkg$ rows displayed by the root packets.

The resulting $\hsame$ comparison is the endpoint comparison required by the
public alpha-beta classifier. Hence the beta-ledger frontier has a deterministic
root endpoint under the displayed classifier.
\end{proof}
